\section{CONCLUSION}\label{ch:conclusion}

I asked whether letting the stock--bond correlation switch with a market regime improves one-day value-at-risk and expected-shortfall forecasts of an equal-weight portfolio of SPY, IEF and GLD. I built a ladder of models that separates the treatment of volatility, correlation and the mean under one regime signal, fixed a single confirmatory test before using 2020--2025, and ran it once.

The test could not tell a gain from none. The estimate favoured switching by 0.0013, the pre-specified test did not reject (outcome O2, Chapter~\ref{ch:holdout}), the sample could not have detected a gain smaller than 0.0031, and a non-significant result is not evidence that the effect is absent.

For a risk manager of a portfolio like this one, the evidence gives no reason to prefer the more complex model, and it does not rule out that the model helps. It also shows a more immediate problem: both primary models breached the 5\% VaR on about 6.8\% to 6.9\% of hold-out days and breach-day losses exceeded the forecast expected shortfall, so neither forecast was well calibrated over 2020--2025. A simple EWMA of the portfolio return, scored post hoc, breached on 5.2\% of days with a loss that cannot be told apart from theirs. On the development data, regime-specific volatility lowered the loss in three of four ETF samples, and the correlation steps were small beside their minimum detectable differences.

The limits are those of Chapter~\ref{ch:discussion}: one portfolio, one vendor, one reversal episode, a regime signal built from equity returns, and a power study that rests on a stylised process. Whether a slower regime signal or a larger sample would show the gain that this one could not is a question for a new, separately pre-specified study.
